%% kapa: exam=1 round=21 year=2010 period=1 form=A subject=economics subject_ko=경제원론 seq=11 number=51
%% source: Q-Net 감정평가사 기출문제 (cst00309) artlSeq=5024181 file=extracted/1008/5024181/2010년 제21회 감정평가사 1차시험 기출문제/1교시.hwp
%% conversion: kordoc 4.24.51 -> tools/gen1.py (auto); review=proofread-sonnet; flags=none
\begin{question}{51}
다음 ( )안의 내용을 옳게 연결한 것은?
\begin{framedtext}
소비함수에 자산효과가 도입되면 물가수준의 하락에 따라 실질자산이 ( ㄱ )하고, 이는 소비의 ( ㄴ )를 통해 ( ㄷ )곡선을 ( ㄹ )으로 이동시켜 국민소득 증가를 가져와 유동성함정 문제를 해결할 수 있다. 이것을 ( ㅁ )효과라고 한다.
\end{framedtext}
\par\smallskip
\hspace{1.2cm}\uline{\makebox[1.6cm]{ㄱ}}\hspace{0.6em}\uline{\makebox[1.6cm]{ㄴ}}\hspace{0.6em}\uline{\makebox[1.6cm]{ㄷ}}\hspace{0.6em}\uline{\makebox[1.6cm]{ㄹ}}\hspace{0.6em}\uline{\makebox[1.6cm]{ㅁ}}\par
\begin{choices}
  \item \makebox[1.6cm]{증가}\hspace{0.6em}\makebox[1.6cm]{증가}\hspace{0.6em}\makebox[1.6cm]{IS}\hspace{0.6em}\makebox[1.6cm]{우측}\hspace{0.6em}\makebox[1.6cm]{케인즈}
  \item \makebox[1.6cm]{증가}\hspace{0.6em}\makebox[1.6cm]{증가}\hspace{0.6em}\makebox[1.6cm]{IS}\hspace{0.6em}\makebox[1.6cm]{우측}\hspace{0.6em}\makebox[1.6cm]{피구}
  \item \makebox[1.6cm]{감소}\hspace{0.6em}\makebox[1.6cm]{감소}\hspace{0.6em}\makebox[1.6cm]{IS}\hspace{0.6em}\makebox[1.6cm]{좌측}\hspace{0.6em}\makebox[1.6cm]{마샬}
  \item \makebox[1.6cm]{증가}\hspace{0.6em}\makebox[1.6cm]{감소}\hspace{0.6em}\makebox[1.6cm]{LM}\hspace{0.6em}\makebox[1.6cm]{좌측}\hspace{0.6em}\makebox[1.6cm]{피구}
  \item \makebox[1.6cm]{감소}\hspace{0.6em}\makebox[1.6cm]{증가}\hspace{0.6em}\makebox[1.6cm]{LM}\hspace{0.6em}\makebox[1.6cm]{우측}\hspace{0.6em}\makebox[1.6cm]{마샬}
\end{choices}
\end{question}
\answer{2}
